\documentclass[10pt,a4paper]{amsart}
\usepackage[margin=19mm]{geometry}
\usepackage{amsmath,amssymb,mathtools}
\usepackage[scr=rsfs,cal=euler]{mathalfa}
\usepackage{xfrac}
\usepackage[unicode,hidelinks,bookmarks=false]{hyperref}
\newcommand{\ie}{i.e.\ }
\newcommand{\cf}{cf.\ }
\newcommand{\resp}{respectively\ }
\newcommand{\Subsection}{\subsection}
\providecommand{\nobreakdash}{\nobreak\hbox{-}}
\setlength{\emergencystretch}{2em}
\multlinegap0pt
\usepackage{bm}
\usepackage{bbm}
\usepackage{dsfont}
\usepackage{xspace}
\usepackage{cancel}
\usepackage{mathtools}
\usepackage{amsthm}
\makeatletter
\@ifundefined{Prop}{\newtheorem{Prop}{Prop}}{}
\@ifundefined{Lem}{\newtheorem{Lem}[Prop]{Lemma}}{}
\makeatother
\title[Two dimensional integral representations via branches of the]{Two dimensional integral representations via branches of the Bruhat-Tits tree}
\author{Bruno Aguil\'o-Vidal, Luis Arenas-Carmona, Mat\'ias Saavedra-Lagos}
\date{Source: arXiv:2506.05661v2 (CC BY 4.0)}
\begin{document}
\maketitle

\noindent\emph{Source and licence.} Two dimensional integral representations via branches of the Bruhat-Tits tree. Version arXiv:2506.05661v2 (CC BY 4.0), distributed under the Creative Commons Attribution 4.0 International licence, as recorded in the arXiv licence metadata for this exact version and shown on the abstract page https://arxiv.org/abs/2506.05661v2. Full rights notice: licenses/arxiv-2506.05661-CC-BY-4.0.txt.\par\smallskip

\noindent\textit{Editorial scope.} Counted unit: the Lem whose source label is \texttt{l61d} in the article (source0.tex lines 2356--2381, source label \texttt{l61d}), delivered together with the article's own complete original proof of it (source0.tex lines 2383--2466, one \texttt{proof} environment of 84 lines). No mathematics is written, repaired or re-derived: statement and proof are the article's own text line for line. Required context retained uncounted: l61c (lines 2290--2307). Every definition, display and label of those lines is retained unchanged. Omitted from this package and not redistributed: 1--2289, 2308--2355, 2382--2382, 2467--3901. Nothing inside a retained excerpt is omitted or altered: every retained body is delivered exactly as the article's lines are written, so no omission receipt of any kind accompanies this package. Citations: the retained excerpts carry no citation command at all, so no bibliography entry of the article is transcribed and no reference is redistributed. Reference adaptation: none was needed and none was made; every reference inside the delivered unit resolves inside it, so no unresolved reference or citation is produced. Printed slips kept literal: no printed word, symbol, hypothesis or number of the counted statement or of its proof was altered. Layout and class: the article's own class is \texttt{article} with packages including babel, geometry, amsmath, amssymb, amsthm, graphicx, enumitem, hyperref, xy, xcolor; the wrapper is the lane's pinned \texttt{amsart} editorial wrapper at 10pt on A4, which restates the theorem-like environments the retained text needs and adds the article's own macro definitions that the retained text needs (none), with extra wrapper packages (bm, bbm, dsfont, xspace, cancel, mathtools). The delivered numbering is the unit's own. Assets: no retained span contains a figure environment, a graphics-inclusion command, a PDF-page inclusion, a table or a TikZ picture; no image file is copied into this package, the package declares no asset dependency and no image file is redistributed with it. Licence: version arXiv:2506.05661v2 of this article is distributed under the Creative Commons Attribution 4.0 International licence, as recorded in the arXiv licence metadata for this exact version and shown on the abstract page https://arxiv.org/abs/2506.05661v2. A full rights notice is delivered at licenses/arxiv-2506.05661-CC-BY-4.0.txt. Characters: every retained excerpt is pure ASCII, so the delivered engine, XeLaTeX with its default Latin Modern text font, reports no missing character.
\begin{Lem}\label{l61c}
     Let $\mathtt{r}$ be a matrix whose eigenvalues $a,b$
     lie in an unramified quadratic extension 
     $L_\omega$ of $K_\nu$,
     and let us denote by $\omega$ the extension 
     of the valuation
     $\nu$ to this field. Then conjugation by $\mathtt{r}$
    fixes precisely the vertices that are at distance at most 
    $\omega\left(1-\frac ab\right)$ from the vertex $v$ 
    in $\mathfrak{t}$ corresponding
    to the unique maximal order in
    $\mathbb{M}_2(K_\nu)$ containing the ring 
    of integers in 
    $\mathfrak{L}_\nu=K_\nu[\mathtt{r}]\cong L_\omega$. 
    The vertex $v$ corresponds to
    a vertex in the stem $\mathfrak{p}$ of the branch
    $\mathfrak{s}_\omega(\mathtt{r})\subseteq \mathfrak{t}'$.
\end{Lem}

\begin{Lem}\label{l61d}
     Let $\mathtt{r}$ be a matrix whose eigenvalues $a,b$
     lie in a ramified quadratic extension 
     $L_\omega$ of $K_\nu$. 
       Let $\omega$ be the (integral valued) 
       extension of $2\nu$
     to the field $L_\omega$. We identify
     $L_\omega$ with the algebra
     $\mathfrak{L}_\nu=K_\nu[\mathtt{r}]$. 
     If $\omega(\mathtt{r})$ is odd,
    then conjugation by $\mathtt{r}$ fixes no vertex
    in the Bruhat-Tits tree $\mathfrak{t}$.
    Furthermore, there is a constant $\kappa$, such that,
    if $\omega(\mathtt{r})$ is even,
    then conjugation by $\mathtt{r}$ 
    fixes precisely the vertices that are at distance at most 
    $d=\frac12\big(\omega
    \left(1-\frac ab\right)-\kappa-1\big)$ 
    from one the vertices 
    corresponding to the two maximal orders 
    containing the ring 
    of integers in 
    $\mathfrak{L}_\nu=K_\nu[\mathtt{r}]$. Furthermore,
    such $d$ is always a positive integer when 
    $\omega(\mathtt{r})$ is even.
\end{Lem}

\begin{proof}
    In this case, the graph that naturally embeds into 
    $\mathfrak{t}'$ is the barycentric subdivision 
    $\mathfrak{t}_b$ of $\mathfrak{t}$, as discussed
    before Lemma \ref{l61c}. It is proved
    as before that the Galois action has a unique invariant
    point in the path $\mathfrak{p}$ connecting the 
    invariant points (visual limits) of
    the Moebius transformation corresponding to
    $\mathtt{r}$, and such invariant
    point is the vertex $v$ that is closest
    to $\mathfrak{t}_b$. In this case, by applying
    a transformation of the form $\mu(z)=az+b$, 
    we can assume that the 
    visual limits of the path are uniformizing parameters 
    of the quadratic extensions, whence the vertex $w$ 
    in $\mathfrak{t}_b$
    that is closest to the path is the barycenter
    of an edge in $\mathfrak{t}$, which, with the 
    assumption on the visual limits, would be the
    edge joining $v_0^{[0]}$ to $v_0^{[1]}$ in
    $\mathfrak{t}$.
    Let $\kappa$ be the distance from $v$ to $w$.

    Let $\mathfrak{H}$ be the ring of integers in 
    $\mathfrak{L}_\nu=K_\nu[\mathtt{r}]
    \subseteq\mathbb{M}_2(K_\nu)$.
    Since $\mathcal{O}_{L_\omega}\mathfrak{H}$ is an order 
    in $L_\omega[\mathtt{r}]$,
    then the maximal orders in 
    $\mathbb{M}_2(L_\omega)$ containing
    $\mathfrak{H}$ correspond to a tubular neighborhood of the 
    path. Therefore, the two maximal orders in 
    $\mathbb{M}_2(K_\nu)$ containing it
    must corresponds to the endpoints of the edge in 
    $\mathfrak{t}$ having $w$ as barycenter. It follows
    that a vertex, in $\mathfrak{t}$, is at distance 
     $\omega\left(1-\frac ab\right)$ from the path,
     if and only if it lies at a distance 
     $\omega\left(1-\frac ab\right)-\kappa$ from $w$, or 
     equivalently,  at a distance 
     $\omega\left(1-\frac ab\right)-\kappa-1$ from one
     of the endpoints of the edge in 
    $\mathfrak{t}$ having a $w$ as barycenter. 
    This is a distance
    in  $\mathfrak{t}'$, so to obtain the corresponding
    distance in  $\mathfrak{t}$ we must divide by two.
    
    Now, fix a uniformizer $\pi_\omega\in L_\omega$, and
    let $\lambda=x+y\pi_\omega\in L_\omega$, 
    with $x,y\in K_\nu$ 
    be arbitrary. Note that $x$ and  
    $y\pi_\omega$ have valuations
    with different parity. If $y\pi_\omega$
    is dominant, then $x/\lambda$
    has positive valuation, whence the relation
    $$\frac{\sigma(\pi_\omega)}{\pi_\omega}-
    \frac{\sigma(\lambda)}{\lambda}=\frac x{\lambda}
    \left(\frac{\sigma(\pi_\omega)}{\pi_\omega}-1\right),$$
    implies that $\frac{\sigma(\lambda)}{\lambda}-1$ 
    has the same 
    valuation as $\frac{\sigma(\pi_\omega)}{\pi_\omega}-1$.
    On the other hand, if $x$ is dominant, then
    $$\frac{\sigma(\lambda)}{\lambda}-1=
    \frac {y\pi_\omega}{\lambda}
    \left(\frac{\sigma(\pi_\omega)}{\pi_\omega}-1\right),$$
    has a larger valuation than 
    $\frac{\sigma(\pi_\omega)}{\pi_\omega}-1$.

    Finally, if $\mathtt{r}$ is a uniformizer 
    in $\mathfrak{L}_\nu$,
    its determinant, as a matrix, is a uniformizer of $K_\nu$, 
    and therefore the Moebius transformation associated to 
    $\mathtt{r}$ cannot leave invariant any vertex of 
    $\mathfrak{t}$. It must, however, leave the path, and 
    therefore the vertex $w$, invariant. It follows that
    $$\omega\Big(\frac{\sigma(\pi_\omega)}{\pi_\omega}-1\Big)=
    \kappa.$$
    Now we conclude from the preceding computations that
    invariant vertices in $\mathfrak{t}$ exists precisely when
    $\omega(\mathtt{r})$ is even, since this coincide with the 
    condition that $\omega\left(1-\frac ab\right)-\kappa-1$
    is non-negative.
\end{proof}

\end{document}
